\chapter{Estado da Arte}
\label{cap:estado-da-arte}

% Apresente aqui uma revisão dos trabalhos e pesquisas mais recentes
% relacionados ao tema, destacando as contribuições e lacunas existentes.

% [Insira o texto do estado da arte aqui]
Esta seção apresenta uma síntese dos trabalhos mais recentes relacionados ao tema, com destaque para suas contribuições e para as lacunas ainda existentes.

\section{Trabalho/Pesquisa A}
\label{sec:trabalho-a}

% [Descreva o primeiro trabalho relacionado]
O primeiro trabalho analisado propõe uma abordagem inovadora para o problema, apresentando resultados promissores em cenários controlados.

\section{Trabalho/Pesquisa B}
\label{sec:trabalho-b}

% [Descreva o segundo trabalho relacionado]
O segundo trabalho, por sua vez, adota uma estratégia distinta, privilegiando a generalização dos resultados. O \autoref{qua:comparativo} sintetiza a comparação entre as duas abordagens.

\begin{quadro}[!htbp]
    \centering
    \Caption{\label{qua:comparativo} Comparação entre as abordagens analisadas}
    \EMBRAPAqua{}{
        \begin{tabular}{|l|c|c|}
            \hline
            \textbf{Critério} & \textbf{Trabalho A} & \textbf{Trabalho B} \\
            \hline
            Abordagem & Inovadora & Tradicional \\
            \hline
            Escopo & Restrito & Amplo \\
            \hline
            Desempenho & Alto & Moderado \\
            \hline
        \end{tabular}
    }{
        \Fonte{Elaborado pelo autor}
    }
\end{quadro}

% Dica: Para inserir um quadro comparativo, use o seguinte modelo:
%
% \begin{quadro}[!htbp]
%     \centering
%     \Caption{\label{qua:modelo} Legenda do quadro}
%     \EMBRAPAqua{}{
%         \begin{tabular}{|c|c|c|}
%             \hline
%             Critério & Trabalho A & Trabalho B \\
%             \hline
%             Critério 1 & Valor & Valor \\
%             \hline
%         \end{tabular}
%     }{
%         \Fonte{Elaborado pelo autor}
%     }
% \end{quadro}
